\chapter{血压调节与运动}\label{ch:b2:blood-pressure}

躺下后迅速站起，约一秒内，有半升
\omterm{def:b2:blood-circulation:blood}{血液}流入腿部的\omterm{def:b2:blood-circulation:vessels}{静脉}；\omterm{def:b2:heart:anatomy}{心脏}接收的\omterm{def:b2:blood-circulation:blood}{血液}减少，
泵出也随之减少，头部压强下降，而大脑——
比\omterm{def:b2:heart:anatomy}{心脏}高 \qty{40}{cm}，且不能储存氧气，
仅一两秒便可能昏厥。这种情况几乎从不发生，因为在一次搏动内，
反射已使\omterm{def:b2:blood-circulation:vessels}{小动脉}收缩、心率加快。跑去赶公交车时，
肌肉所需血量是静息时的十倍；
它们可在一分钟内获得\omterm{def:b2:blood-circulation:blood}{血液}，驱动血流的压强
保持稳定，而全身阻力降低
三分之二。本章讨论这一调节：何者决定血压，
血压如何被感知，快速反射与缓慢激素如何纠正血压，
以及运动时该系统如何被重组，而非
被压制。

\section{血压的含义}

\begin{theorem}[循环方程]\label{thm:b2:blood-pressure:equation}
\index{平均动脉压}\index{总外周阻力}
\omterm{thm:b2:blood-pressure:equation}{平均动脉压} $\bar P$ 是\omterm{def:b2:heart:anatomy}{心脏}泵入量
与血管排出量的乘积：
\[
\bar P - P_{\text{静脉}} = Q\,R, \qquad Q = f\times V_s,
\]
其中 $Q$ 为\omterm{thm:b2:heart:starling}{心输出量}（心率 $f$ 乘以\omterm{prop:b2:heart:cycle}{每搏输出量}
$V_s$），$R$ 为\omterm{thm:b2:blood-pressure:equation}{总外周阻力}——所有器官的\omterm{def:b2:blood-circulation:vessels}{小动脉}
并联阻力，$P_{\text{静脉}}$ 为\omterm{def:b2:blood-circulation:vessels}{静脉}压，
接近零。静息时，$\qty{5}{L/min}\times\qty{1.08}{mmHg.s/mL}$ 给出
\qty{90}{mmHg}。血压的任何调节，都作用于下列三个
量之一：心率、\omterm{prop:b2:heart:cycle}{每搏输出量}（通过充盈和
收缩性（见\cref{ch:b2:heart}）或阻力（通过
\omterm{def:b2:blood-circulation:vessels}{小动脉}半径，按四次方关系）。各器官并联，
故流向每个器官的血流等于其所分得的压强除以
自身阻力：扩张某一器官的\omterm{def:b2:blood-circulation:vessels}{小动脉}可使该器官血流增加，
而不改变其他器官的血流，前提是压强保持不变——
维持压强正是调节的目的。
\end{theorem}
\begin{proof}
流体的欧姆定律：流经阻力的流量
等于该阻力两端的压降除以阻力（泊肃叶，
见\cref{ch:b2:blood-circulation}）；并联阻力时，
流量相加，且 $1/R = \sum 1/R_i$。一次\omterm{prop:b2:heart:cycle}{心动周期}内的平均压强
约为舒张压加上\omterm{thm:b2:blood-circulation:windkessel}{脉压}的三分之一：
$80 + 40/3 = \qty{93}{mmHg}$。
\end{proof}

\begin{omfigure}
\begin{tikzpicture}[every node/.style={font=\scriptsize}, >=stealth]
  \node[draw, rounded corners, fill=omThm!15, minimum width=2.4cm, minimum height=1cm, align=center] (heart) at (0,0) {心脏 \\$Q = f\,V_s$};
  \draw[very thick, omThm] (1.2,0) -- (2.4,0);
  \draw[very thick, omThm] (2.4,-2.4) -- (2.4,2.4);
  \foreach \y/\org/\r in {2.0/{脑}/{$R_1$}, 1.0/{心肌}/{$R_2$}, 0/{肠和肝}/{$R_3$}, -1.0/{肾}/{$R_4$}, -2.0/{肌肉、皮肤}/{$R_5$}} {
    \draw[very thick, omThm] (2.4,\y) -- (3.6,\y);
    \draw[thick, decorate, decoration={zigzag, segment length=4pt, amplitude=2pt}] (3.6,\y) -- (5.0,\y);
    \node[above] at (4.3,\y+0.05) {\org};
    \node[below] at (4.3,\y-0.05) {\r};
    \draw[very thick, omDef] (5.0,\y) -- (6.2,\y);
  }
  \draw[very thick, omDef] (6.2,-2.4) -- (6.2,2.4);
  \draw[very thick, omDef] (6.2,0) -- (7.4,0) |- (0,-3.2) -- (0,-0.5);
  \node[omThm, above] at (1.8,0.05) {$\bar P$};
  \node[omDef, below] at (3.7,-3.2) {静脉，$P_{\text{静脉}} \approx 0$};
  \node[anchor=west, align=left] at (7.8,1.5) {并联器官：\\$1/R = \sum 1/R_i$\\流向器官 $i$ 的血流：$\bar P/R_i$};
\end{tikzpicture}

{\small 循环系统作为一个回路：一个泵、一个压强，以及
各器官并联的\omterm{def:b2:blood-circulation:vessels}{小动脉}阻力。每个器官获得
$\bar P/R_i$ 的血流；调节使 $\bar P$ 保持不变，同时 $R_i$ 随
需求而调整。}
\end{omfigure}

\begin{proposition}[为何必须调节血压]\label{prop:b2:blood-pressure:why}
压强过低时，大脑——站立时比\omterm{def:b2:heart:anatomy}{心脏}高 \qty{40}{cm}，
成年人需克服 \qty{31}{mmHg} 的静水压头——以及
心肌本身灌注不足：大脑在
数秒内即发生功能衰竭。压强过高，则血管受损：\omterm{def:b2:blood-circulation:vessels}{动脉}壁
增厚变硬，肾小球瘢痕化，视网膜
出血，\omterm{def:b2:heart:anatomy}{心脏}则因对抗高压而劳损，直至衰竭。
此外，在各器官需求变化达一个数量级时，压强也必须保持稳定——
例如餐后的肠道、运动时的肌肉、
高温时的皮肤——同时血容量本身
因出血或出汗而减少。两套系统承担此任务：一套是神经
反射，在一次搏动内作用于心率、\omterm{prop:b2:heart:cycle}{每搏输出量}和
阻力；另一套是激素系统，在数小时至数日内
作用于阻力，并且最主要通过
肾脏作用于血容量。
\end{proposition}

\begin{omfigure}
\includegraphics[width=0.6\linewidth]{images/book4/ai/b4-18-blood-pressure-cuff.jpg}

{\small 测量血压。袖带充气至高于
收缩压，使血流不通；放气时，
听诊器下最先出现的湍流声标志
收缩压，这些声音消失时，\omterm{def:b2:blood-circulation:vessels}{动脉}
在整个周期中保持开放，标志舒张压。}
\end{omfigure}

\section{快速回路：压力反射}

\begin{definition}[压力反射]\label{def:b2:blood-pressure:baroreflex}
\index{压力感受器}\index{压力反射}\index{血管运动中枢}\index{自主神经系统}
牵张感受器位于颈\omterm{def:b2:blood-circulation:vessels}{动脉}窦壁（每条颈\omterm{def:b2:blood-circulation:vessels}{动脉}的分叉处，
位于下颌下方）以及主\omterm{def:b2:blood-circulation:vessels}{动脉}弓的壁内——即
\emph{压力感受器}——放电频率随
牵张它们的压强升高而增加，压强上升时更快。它们的神经
投射至延髓的\emph{心血管中枢}，该中枢
控制自主神经传出的两个分支：迷走神经，其
乙酰胆碱使\omterm{prop:b2:heart:conduction}{窦房结}减慢；交感神经，
其去甲肾上腺素加快\omterm{prop:b2:heart:conduction}{窦房结}、增强
\omterm{def:b2:heart:anatomy}{心室}收缩力、使\omterm{def:b2:blood-circulation:vessels}{小动脉}收缩（使 $R$ 增大），并使
\omterm{def:b2:blood-circulation:vessels}{静脉}收缩（将贮存\omterm{def:b2:blood-circulation:blood}{血液}挤向\omterm{def:b2:heart:anatomy}{心脏}，提高
充盈量）。压强升高时压力感受器放电增加，此信号
兴奋迷走神经，抑制交感传出：心率
下降，\omterm{def:b2:blood-circulation:vessels}{小动脉}舒张，压强下降。压强下降则
作用相反。该环路是一个\emph{负反馈}，其调定点
约为 \qty{95}{mmHg}，延迟约一秒，增益
使扰动被纠正至其幅度的五分之一以内：
这正是你站立时不会昏厥的原因。
\end{definition}
\begin{proof}[证据]
赫林（1923 年）发现，按压颈\omterm{def:b2:blood-circulation:vessels}{动脉}窦可使
心率减慢、血压下降；切断窦神经则
该反应消失。在设定压强下灌注离体颈\omterm{def:b2:blood-circulation:vessels}{动脉}窦，
同时记录动物的\omterm{def:b2:blood-circulation:vessels}{动脉}压，得到
该反射的S形曲线：随着窦内压升高，动物血压
下降；在 \qty{100}{mmHg} 附近陡降，
而在 \qtyrange{60}{160}{mmHg} 之外变化很小。同时切断窦神经与主\omterm{def:b2:blood-circulation:vessels}{动脉}神经的犬
保持数日的正常平均压，
但血压逐分钟剧烈波动——该反射是缓冲器，
而非长期水平的设定者。
\end{proof}

\begin{omfigure}
\begin{tikzpicture}[every node/.style={font=\scriptsize}, >=stealth]
  \node[draw, rounded corners, fill=omDef!15, align=center] (bp) at (0,0) {动脉压 \\$\bar P$};
  \node[draw, rounded corners, fill=omDef!15, align=center] (baro) at (4.2,0) {压力感受器 \\颈动脉窦、主动脉弓 \\放电频率 $\uparrow$，随 $\bar P$ 升高};
  \node[draw, rounded corners, fill=omProp!20, align=center] (med) at (8.6,0) {延髓：\\心血管中枢};
  \node[draw, rounded corners, fill=omThm!15, align=center] (vag) at (8.6,-2.2) {迷走神经 $\uparrow$：心率 $\downarrow$};
  \node[draw, rounded corners, fill=omThm!15, align=center] (sym) at (4.2,-2.2) {交感神经 $\downarrow$：心率、收缩力，\\小动脉和静脉张力 $\downarrow$};
  \draw[->, thick] (bp) -- (baro); \draw[->, thick] (baro) -- (med);
  \draw[->, thick] (med) -- (vag); \draw[->, thick] (med) -- (sym);
  \draw[->, thick] (sym.west) -| (bp.south) node[pos=0.75, left, align=right] {$Q\downarrow$, $R\downarrow$:\\$\bar P\downarrow$};
  \draw[->, thick] (vag.west) -- (sym.east);
  \node[anchor=north, align=center] at (4.3,-3.2) {压强升高则随之下降：负反馈，延迟一秒};
\end{tikzpicture}

{\small \omterm{def:b2:blood-pressure:baroreflex}{压力反射}。\omterm{def:b2:blood-circulation:vessels}{动脉}压升高，使
\omterm{def:b2:blood-pressure:baroreflex}{压力感受器}放电增加，驱动延髓减慢心率并
舒张血管；压力下降则作用相反。}
\end{omfigure}

\begin{theorem}[负反馈回路的增益]\label{thm:b2:blood-pressure:gain}
\index{反馈增益}
若某一扰动在无调节时使压强变化
$\Delta P_0$，而回路对任一偏差 $\Delta P$ 的
响应为 $-G\,\Delta P$ 的校正（\emph{开环增益} $G$），
则回路作用后残留的偏差为
\[
\Delta P = \frac{\Delta P_0}{1 + G}.
\]
\omterm{def:b2:blood-pressure:baroreflex}{压力反射}的 $G \approx 4$：站起时，原本
会使\omterm{def:b2:heart:anatomy}{心脏}处的压强下降约 \qty{40}{mmHg}，实际仅下降 \qty{8}{mmHg}；
若失血会使压强减半，反射则只使其降低
十分之一——直到失血量超过血管收缩与心跳加快所能掩盖的限度。
反射并不能消除扰动，
而是把它缩小为原来的五分之一。此外它会\emph{重设}：
在新的压强下维持数小时，\omterm{def:b2:blood-pressure:baroreflex}{压力感受器}便适应并维护新的水平，这
正是该反射无法根治高血压的原因，也说明为何
长期水平由别处设定。
\end{theorem}
\begin{proof}
终态偏差等于扰动加上校正：
$\Delta P = \Delta P_0 - G\,\Delta P$，于是 $\Delta P(1 + G) =
\Delta P_0$。取 $G = 4$，则 $\Delta P = \Delta P_0/5$。
\end{proof}

\begin{omfigure}
\begin{tikzpicture}
\begin{axis}[
  width=0.7\linewidth, height=5.4cm,
  axis x line=bottom, axis y line=left,
  xmin=40, xmax=180, ymin=40, ymax=160,
  xlabel={颈动脉窦内压强（\unit{mmHg}）}, ylabel={随之达到的动脉压（\unit{mmHg}）},
  xlabel style={font=\small, at={(axis description cs:0.5,-0.13)}, anchor=north},
  ylabel style={font=\small},
  tick label style={font=\small}, xtick={60,100,140}, ytick={60,100,140},
  clip=false,
]
\addplot[very thick, omDef, domain=40:180, samples=80] {100 + 50*tanh((100-x)/25)};
\addplot[thick, dashed, gray, domain=40:180, samples=2] {x};
\fill[omDef] (axis cs:100,100) circle (2pt); \draw[gray] (axis cs:101,100) -- (axis cs:124,110); \node[font=\scriptsize, anchor=west] at (axis cs:125,110) {调定点};
\node[font=\scriptsize, anchor=west, align=left] at (axis cs:46,112) {斜率 $-G$\\调定点附近};
\end{axis}
\end{tikzpicture}

{\small \omterm{def:b2:blood-pressure:baroreflex}{压力反射}曲线：分离的颈\omterm{def:b2:blood-circulation:vessels}{动脉}窦维持在给定压强时，机体
最终稳定的\omterm{def:b2:blood-circulation:vessels}{动脉}压。曲线在调定点处的斜率即为增益；
在 \qtyrange{60}{160}{mmHg} 之外，
反射已饱和。}
\end{omfigure}

\section{慢回路：激素与肾脏}

\begin{proposition}[肾素、血管紧张素、醛固酮与血容量]\label{prop:b2:blood-pressure:raas}
\index{肾素–血管紧张素系统}\index{醛固酮}\index{抗利尿激素}\index{利钠肽}
在数小时至数天的时间尺度上，压强由\emph{血容量}决定，
而血容量则由肾脏调控。压强下降、或输送至肾脏的钠减少，
或交感神经放电时，肾\omterm{def:b2:blood-circulation:vessels}{小动脉}的细胞
将酶\emph{肾素}释放入
\omterm{def:b2:blood-circulation:blood}{血液}；肾素将\omterm{def:b2:blood-circulation:blood}{血浆}蛋白裂解为\emph{血管紧张素 I}，
肺\omterm{def:b2:blood-circulation:vessels}{毛细血管}中的一种酶将其转化为\emph{血管紧张素 II}
——体内最强的缩血管物质，
可在数分钟内使 $R$ 升高；它还是一种激素，促使肾上腺皮质分泌
\emph{\omterm{prop:b2:blood-pressure:raas}{醛固酮}}，使肾脏保留钠，并随之
保留水分，历时数日。垂体分泌的
\emph{\omterm{prop:b2:blood-pressure:raas}{抗利尿激素}}（血管加压素）于\omterm{def:b2:blood-circulation:blood}{血液}浓缩或血容量下降时释放，
可使肾脏直接保留水分，同时收缩血管。当\omterm{def:b2:heart:anatomy}{心房}
因容量过多而过度扩张时，
便分泌\emph{\omterm{prop:b2:blood-pressure:raas}{利钠肽}}，
促使肾脏排出钠和水。
因此肾脏是最终仲裁者：它可排出或保留数升，
而持续高于肾脏调定点的压强，会
经数日通过血容量减少而被纠正——除非肾脏的
调定点本身被提高，这正是多数高血压的情形。
（肾脏自身的机制是本科三年级分册的内容。）
\end{proposition}
\begin{proof}[证据]
戈德布拉特（1934 年）收窄了通向一只犬某一侧肾脏的\omterm{def:b2:blood-circulation:vessels}{动脉}：
该动物数日内即患高血压，并持续不退；
灌注不足的肾脏被发现会释放肾素，切除该肾
即可治愈该动物。阻断血管紧张素转化的药物
（ACE 抑制剂）或阻断其受体的药物，可降低多数
高血压患者的血压；将高血压大鼠品系的肾脏
移植到正常受体身上，可使其患高血压。
\end{proof}

\begin{omfigure}
\begin{tikzpicture}[every node/.style={font=\scriptsize}, >=stealth]
  \node[draw, rounded corners, fill=omDef!15, align=center] (fall) at (0,0) {压强或血容量下降；\\交感神经放电};
  \node[draw, rounded corners, fill=omProp!20, align=center] (renin) at (4.4,0) {肾脏释放 \\肾素};
  \node[draw, rounded corners, fill=omProp!20, align=center] (ang) at (8.6,0) {血管紧张素 II\\（在血液和肺中生成）};
  \node[draw, rounded corners, fill=omThm!15, align=center] (vc) at (8.6,-2.0) {血管收缩：\\$R\uparrow$，数分钟内};
  \node[draw, rounded corners, fill=omThm!15, align=center] (ald) at (4.4,-2.0) {醛固酮（肾上腺）：\\Na$^+$ 与水被保留 \\数日内——血容量 $\uparrow$};
  \draw[->, thick] (fall) -- (renin); \draw[->, thick] (renin) -- (ang);
  \draw[->, thick] (ang) -- (vc); \draw[->, thick] (ang) -- (ald);
  \draw[->, thick] (ald.west) -| (fall.south) node[pos=0.85, left, align=right] {压强 \\恢复};
  \draw[->, thick] (vc.west) -- (ald.east);
  \node[draw, rounded corners, fill=omExo!25, align=center] (anp) at (0,-2.0) {心房过度扩张：\\利钠肽，\\钠和水排出};
  \draw[->, thick, dashed] (anp.north) -- (fall.south);
\end{tikzpicture}

{\small 慢回路。血压下降引发肾素、
血管紧张素 II 与\omterm{prop:b2:blood-pressure:raas}{醛固酮}：数分钟内血管收缩，盐与水
潴留则历时数日。血容量过多时，
由\omterm{def:b2:heart:anatomy}{心房}分泌的\omterm{prop:b2:blood-pressure:raas}{利钠肽}加以应对。}
\end{omfigure}

\begin{example}[失血]\label{ex:b2:blood-pressure:haemorrhage}
一名施主失血 \qty{500}{mL}，占\omterm{def:b2:blood-circulation:blood}{血液}的十分之一。数秒之内，
\omterm{def:b2:blood-pressure:baroreflex}{压力反射}使\omterm{def:b2:blood-circulation:vessels}{小动脉}和\omterm{def:b2:blood-circulation:vessels}{静脉}收缩，并加快
心率：血压几乎不变，皮肤变得苍白、发凉，
（其\omterm{def:b2:blood-circulation:vessels}{小动脉}已关闭），脉搏加快。数分钟之内，
\omterm{def:b2:blood-circulation:vessels}{毛细血管}压力降低，组织液进入血管
（即\cref{ch:b2:blood-circulation}中的斯塔林平衡），
约一小时内恢复一半容量；肾素、血管紧张素与
\omterm{prop:b2:blood-pressure:raas}{抗利尿激素}使尿量减至涓滴，口渴
随之而来；数日内肾脏保留盐与水，\omterm{def:b2:blood-circulation:blood}{血浆}
容量恢复，数周内骨髓补充\omterm{def:b2:blood-circulation:blood}{红细胞}。若一次
失血达 \qty{2}{L}，则超出上述全部代偿：\omterm{def:b2:blood-pressure:baroreflex}{压力反射}已饱和，
血压骤降，灌注不足的组织释放乳酸，
\omterm{def:b2:blood-circulation:vessels}{毛细血管}渗漏——休克。只有输血
才能使患者恢复。
\end{example}

\section{练习：系统的重新组织}

\begin{proposition}[练习]\label{prop:b2:blood-pressure:exercise}
\index{运动生理学}\index{肌肉血流}\index{心输出量!运动}
剧烈运动时，骨骼肌的耗氧量上升
为原来的二十倍，其血流量由 \qty{1}{L/min} 增至 \qty{20}{L/min}。
运动时同时发生三件事。局部上，工作肌肉的代谢产物——
$\mathrm{CO_2}$、乳酸、腺苷、钾、热量、
低氧——使其\omterm{def:b2:blood-circulation:vessels}{小动脉}舒张（\emph{代谢性血管舒张}），
并使\omterm{def:b2:blood-circulation:vessels}{毛细血管}开放，将该肌肉的阻力降至
原来的十分之一。中枢方面，运动皮层发往延髓的指令，
加上肌肉自身受体发出的信号予以增强，将
\omterm{def:b2:blood-pressure:baroreflex}{压力反射}的调定点重设于更高水平，并驱动交感神经输出：
心率升至 180，收缩力增强，肠道、肾脏与
静息肌肉的\omterm{def:b2:blood-circulation:vessels}{小动脉}收缩，\omterm{def:b2:blood-circulation:vessels}{静脉}受挤压；肌泵与
深呼吸共同加快\omterm{def:b2:blood-circulation:blood}{血液}回流。结果是
输出量达 \qty{25}{L/min}，总阻力为静息值的三分之一，
收缩压为 \qty{150}{mmHg}，舒张压则不高于静息时；
五分之四的血流流向肌肉，脑的份额
按绝对量计并未改变，肠道的份额减半，
皮肤的份额增加以散热。\omterm{def:b2:blood-pressure:baroreflex}{压力反射}并非被推翻，
而是被重新调校：它维持较高的血压，同时使
全身血管重新开放。
\end{proposition}

\begin{omfigure}
\begin{tikzpicture}
\begin{axis}[
  ybar, bar width=9pt,
  width=0.9\linewidth, height=6cm,
  ymin=0, ymax=22, ytick={0,5,10,15,20},
  ylabel={血流量（\unit{L/min}）}, ylabel style={font=\small},
  symbolic x coords={brain, heart, gut and liver, kidneys, muscle, skin, other},
  xtick=data, xticklabels={脑, 心脏, 肠和肝, 肾, 肌肉, 皮肤, 其他}, x tick label style={font=\scriptsize, rotate=30, anchor=north east},
  tick label style={font=\small},
  legend style={font=\scriptsize, at={(0.03,0.97)}, anchor=north west, draw=none},
  enlarge x limits=0.12,
  clip=false,
]
\addplot[fill=omDef!50, draw=omDef] coordinates {(brain,0.75) (heart,0.25) (gut and liver,1.4) (kidneys,1.1) (muscle,1.0) (skin,0.3) (other,0.2)};
\addlegendentry{静息，\qty{5}{L/min}}
\addplot[fill=omThm!50, draw=omThm] coordinates {(brain,0.75) (heart,1.0) (gut and liver,0.6) (kidneys,0.6) (muscle,20) (skin,1.6) (other,0.45)};
\addlegendentry{剧烈运动，\qty{25}{L/min}}
\end{axis}
\end{tikzpicture}

{\small \omterm{def:b2:blood-circulation:blood}{血液}的去向。剧烈运动时，输出量
升至五倍，并重新分配：流向肌肉的增至二十倍，流向
\omterm{def:b2:heart:anatomy}{心脏}与皮肤的增多，流向肠道与肾脏的减少，流向
脑的则不变。}
\end{omfigure}

\begin{theorem}[氧的输送与菲克原理]\label{thm:b2:blood-pressure:fick}
\index{菲克原理}\index{VO2max@$\dot V_{\mathrm{O_2}\max}$}
机体每分钟消耗的氧量，等于\omterm{thm:b2:heart:starling}{心输出量}
乘以\omterm{def:b2:blood-circulation:vessels}{动脉}血与混合
\omterm{def:b2:blood-circulation:vessels}{静脉}血的氧含量之差：
\[
\dot V_{\mathrm{O_2}} = Q\,(C_a - C_v).
\]
静息时，$\qty{5}{L/min}\times(200 - 150)\,\unit{mL/L} = \qty{250}{mL/min}$；
最大运动时，$\qty{25}{L/min}\times(200 - 40) =
\qty{4000}{mL/min}$，增至十六倍，这源于血流增至五倍
以及摄取率增至三倍。这一最大速率，
$\dot V_{\mathrm{O_2}\max}$，是衡量耐力
适合度的标准指标，其上限取决于\omterm{def:b2:heart:anatomy}{心脏}：肌肉本可再多摄取氧，
肺也可再多载氧，但输出量已无法
再升。训练可提高\omterm{prop:b2:heart:cycle}{每搏输出量}（\omterm{def:b2:heart:anatomy}{心室}更大、更具顺应性，
血容量更多），从而提高输出量——
运动员静息心率为 45，正是\omterm{prop:b2:heart:cycle}{每搏输出量}
达 \qty{110}{mL} 的标志——并为肌肉增加\omterm{def:b2:blood-circulation:vessels}{毛细血管}与线粒体，
使其摄取更多。反之，菲克原理是测定患者
\omterm{thm:b2:heart:starling}{心输出量}的测定方法：依据所耗氧量与
两份血样。
\end{theorem}
\begin{proof}
每升流经组织的\omterm{def:b2:blood-circulation:blood}{血液}留下 $C_a - C_v$
毫升氧，而每分钟流过 $Q$ 升\omterm{def:b2:blood-circulation:blood}{血液}；稳态时，
这即是肺摄取并由机体消耗的量。
由此解出 $Q$，即得测量值。
\end{proof}

\begin{example}[起立、晕厥与宇航员]\label{ex:b2:blood-pressure:standing}
起立时，\qty{500}{mL}\omterm{def:b2:blood-circulation:blood}{血液}转入腿部\omterm{def:b2:blood-circulation:vessels}{静脉}，
回心血量与\omterm{prop:b2:heart:cycle}{每搏输出量}下降三分之一，颈\omterm{def:b2:blood-circulation:vessels}{动脉}窦
处压力降低——\omterm{def:b2:blood-pressure:baroreflex}{压力反射}在一次搏动内作出应答，
心率加快、血管收紧，头部压力
在十秒之内恢复。一名士兵在炎热的检阅场上
静立不动，失去肌泵作用，\omterm{def:b2:blood-circulation:blood}{血液}淤积于扩张的皮肤\omterm{def:b2:blood-circulation:vessels}{静脉}，
于是昏厥：平躺即刻使其恢复。一位
宇航员在数月失重之后，血容量减少了一升
（\omterm{def:b2:blood-circulation:blood}{血液}淤积于胸腔时，肾脏将其误读为过剩而排出），
而着陆后，久未练习的反射无法
维持其直立。长颈鹿的头部高出\omterm{def:b2:heart:anatomy}{心脏}两米，
其平均血压为 \qty{200}{mmHg}，并穿着
由皮肤构成的加压袜——
\cref{ch:b2:blood-circulation}中的流体静力学提出了要求，
而本章的反射则满足这些要求。
\end{example}

\section{练习}

\begin{exercise}[$\star$]\label{exo:b2:blood-pressure:1}
写出平均压、\omterm{thm:b2:heart:starling}{心输出量}与外周阻力之间的关系式，
并列出机体能够改变的三个量，
以及每个量发生变化的器官或组织。
\end{exercise}

\begin{exercise}[$\star$]\label{exo:b2:blood-pressure:2}
描述\omterm{def:b2:blood-pressure:baroreflex}{压力反射}：感受器、中枢、效应器、
反馈的符号与延迟。若将某动物的\omterm{def:b2:blood-pressure:baroreflex}{压力感受器}神经切断，
其压强逐分钟如何变化？
\end{exercise}

\begin{exercise}[$\star$]\label{exo:b2:blood-pressure:3}
给出肾素 $\to$ 血管紧张素 $\to$ \omterm{prop:b2:blood-pressure:raas}{醛固酮}的序列，
说明各步由哪一器官产生、各自起何作用，
并指出是什么启动了这一过程。
\end{exercise}

\begin{exercise}[$\star$]\label{exo:b2:blood-pressure:4}
指出使工作\omterm{prop:b2:blood-pressure:exercise}{肌肉血流}量增加二十倍的三种机制，
并说明为此哪些器官需要让出血流，
以使之成为可能。
\end{exercise}

\begin{exercise}[$\star\star$]\label{exo:b2:blood-pressure:5}
计算 $Q = \qty{5}{L/min}$ 与 $R =
\qty{1.2}{mmHg.s/mL}$ 时的平均压，以及 $Q = \qty{20}{L/min}$ 与 $R =
\qty{0.35}{mmHg.s/mL}$ 时的平均压。在第二种情形中，平均\omterm{def:b2:blood-circulation:vessels}{小动脉}半径
变为原来的多少倍？
\end{exercise}

\begin{exercise}[$\star\star$]\label{exo:b2:blood-pressure:6}
一次扰动会使压强下降 \qty{30}{mmHg}。计算
反射增益为 2、4 和 8 时的剩余降幅。某药物使患者的\omterm{def:b2:blood-pressure:baroreflex}{压力反射}增益减半，
她站立时会有何感觉？
\end{exercise}

\begin{exercise}[$\star\star$]\label{exo:b2:blood-pressure:7}
\omterm{def:b2:blood-circulation:blood}{血液}密度为 \qty{1060}{kg/m^3}；站立的成年人脑部比\omterm{def:b2:heart:anatomy}{心脏}高 \qty{40}{cm}，
足部比\omterm{def:b2:heart:anatomy}{心脏}低 \qty{120}{cm}。计算
各处的流体静压差（以 mmHg 计），以及
当\omterm{def:b2:heart:anatomy}{心脏}处\omterm{def:b2:blood-circulation:vessels}{动脉}压为
\qty{95}{mmHg} 时，脑部与踝部的\omterm{def:b2:blood-circulation:vessels}{动脉}压。一头头部比\omterm{def:b2:heart:anatomy}{心脏}高 \qty{2.5}{m} 的长颈鹿
需要什么？
\end{exercise}

\begin{exercise}[$\star\star$]\label{exo:b2:blood-pressure:8}
利用菲克原理计算一位患者的\omterm{thm:b2:heart:starling}{心输出量}，其每分钟耗氧
\qty{280}{mL}，\omterm{def:b2:blood-circulation:vessels}{动脉}血氧含量为 \qty{190}{mL/L}，混合\omterm{def:b2:blood-circulation:vessels}{静脉}血氧含量为
\qty{130}{mL/L}。一名运动员在最大强度下耗氧
\qty{5}{L/min}，\omterm{def:b2:blood-circulation:vessels}{动脉}与混合\omterm{def:b2:blood-circulation:vessels}{静脉}的氧含量分别为 200 和
\qty{30}{mL/L}，计算其\omterm{thm:b2:heart:starling}{心输出量}，并在心率为 190 时计算其\omterm{prop:b2:heart:cycle}{每搏
输出量}。
\end{exercise}

\begin{exercise}[$\star\star$]\label{exo:b2:blood-pressure:9}
一次 \qty{800}{mL} 的出血若无调节，会使压强下降
\qty{35}{mmHg}。给出\omterm{def:b2:blood-pressure:baroreflex}{压力反射}（$G = 4$）作用后的压强，
列出恢复血容量的四种机制及其时间尺度，
并解释皮肤苍白、发凉的原因。
\end{exercise}

\begin{exercise}[$\star\star\star$]\label{exo:b2:blood-pressure:10}
静息时肌肉的 $R = \qty{5}{mmHg.s/mL}$，接受的血流为
\qty{1}{L/min}；运动时其阻力降至十分之一。若压强保持在
\qty{95}{mmHg}，其他条件不变，
肌肉将获得多少血流？\omterm{def:b2:heart:anatomy}{心脏}需要输出多少？
若改为\omterm{thm:b2:heart:starling}{心输出量}固定为 \qty{5}{L/min}，
压强将降至多少？解释机体为何两者都不做，
以及它实际做了什么。
\end{exercise}

\begin{exercise}[$\star\star\star$]\label{exo:b2:blood-pressure:11}
戈德布拉特缩窄了一侧肾\omterm{def:b2:blood-circulation:vessels}{动脉}，犬随之出现持续性高血压，
而\omterm{def:b2:blood-pressure:baroreflex}{压力反射}的反应仍然正常。解释 (a) 为何肾脏会升高血压，
(b) 为何\omterm{def:b2:blood-pressure:baroreflex}{压力反射}未能阻止这一升高，
(c) 为何切除该肾可治愈高血压，
(d) 若使用 ACE 抑制剂，会产生什么效果。
\end{exercise}

\begin{exercise}[$\star\star\star$]\label{exo:b2:blood-pressure:12}
“\omterm{def:b2:blood-pressure:baroreflex}{压力反射}是缓冲器，肾脏是恒温器。”
讨论这两个系统的时间尺度、增益与调定点，
以及两者各自能够、不能够对慢性升高的血压做什么。
\end{exercise}

\section{问题：站立、出血、奔跑}

\begin{problem}[{周末问题——一套循环系统经历三项挑战，
并计算其压强、血流、反射校正与氧输送，
最后落在站立时的剩余降幅、
出血后的压强，以及
运动中的心输出量与摄氧量}]
\label{pb:b2:blood-pressure:1}
静息数据：$Q = \qty{5}{L/min}$，心率 70，$\bar P = \qty{95}{mmHg}$，
$P_{\text{静脉}} = 0$，\omterm{def:b2:blood-circulation:blood}{血液} \qty{5}{L}，\omterm{def:b2:blood-pressure:baroreflex}{压力反射}增益 $G = 4$，
\omterm{def:b2:blood-circulation:vessels}{动脉}血氧 \qty{200}{mL/L}，\omterm{def:b2:blood-circulation:vessels}{静脉}血氧 \qty{150}{mL/L}。器官
静息阻力（\unit{mmHg.s/mL}）：脑 7.6，\omterm{def:b2:heart:anatomy}{心脏} 22.8，肠道
与肝脏 4.1，肾脏 5.2，肌肉 5.7，皮肤 19，其他 28.5。\omterm{def:b2:blood-circulation:blood}{血液}
密度为 \qty{1060}{kg/m^3}；$\qty{1}{mmHg} = \qty{133}{Pa}$。

\medskip
\noindent\textbf{第一部分——静息状态。}
\begin{enumerate}
  \item 由各器官阻力并联计算\omterm{thm:b2:blood-pressure:equation}{总外周阻力}，
        并与 $\bar P/Q$ 作对照检验。
  \item 计算流向各器官的血流，
        以及各器官所获得的占\omterm{thm:b2:heart:starling}{心输出量}的比例。
  \item 计算\omterm{prop:b2:heart:cycle}{每搏输出量}与耗氧量。
  \item 脑重 \qty{1.4}{kg}，肾脏重 \qty{0.3}{kg}。
        计算它们每千克的血流，并加以评论。
  \item 为何必须以压强而非血流作为
        受调节的变量？
  \item 心肌每分钟接受 \qty{250}{mL/min}，提取其中
        \qty{70}{\%} 的氧；静息时的肌肉提取
        \qty{25}{\%}。计算两者每分钟所用的氧量，
        并解释为何心肌唯一的储备在于增加血流。
\end{enumerate}

\medskip
\noindent\textbf{第二部分——站起身。}
\begin{enumerate}[resume]
  \item 站起时，\qty{500}{mL} 转入腿部\omterm{def:b2:blood-circulation:vessels}{静脉}，使
        \omterm{prop:b2:heart:cycle}{每搏输出量}下降 \qty{30}{\%}。计算
        无调节时 $\bar P$ 的下降量（阻力不变）。
  \item 计算反射作用后的剩余下降量。
  \item 反射通过将心率升至 85 并提高
        阻力来实现此作用。计算所需的新阻力。
  \item 计算\omterm{def:b2:heart:anatomy}{心脏}与比其高 \qty{40}{cm} 的脑部之间的
        静水压差，以及反射前后该脑部的
        \omterm{def:b2:blood-circulation:vessels}{动脉}压。
  \item 一名服用阻断交感神经药物的患者，其
        $G = 1$。计算她剩余的下降量与脑部压强，
        并说明她会有何感受。
  \item 解释为何将昏厥的患者平躺，
        可在数秒内恢复意识。
\end{enumerate}

\medskip
\noindent\textbf{第三部分——失血。}
\begin{enumerate}[resume]
  \item 若失血 \qty{1}{L}，无调节时 $\bar P$ 将降低
        \qty{45}{mmHg}。计算反射作用后的压强。
  \item 反射使皮肤与肠道的\omterm{def:b2:blood-circulation:vessels}{小动脉}收缩，使其
        阻力加倍。取第 13 题的压强，重新计算
        总阻力，以及流向皮肤、肠道与脑部的
        血流。
  \item 在接下来的一小时内，\qty{500}{mL} 组织液
        进入\omterm{def:b2:blood-circulation:vessels}{毛细血管}。用\omterm{prop:b2:blood-circulation:exchange}{斯塔林力}解释其机制，
        以及\omterm{def:b2:blood-circulation:blood}{血细胞比容}会发生怎样的变化。
  \item 三天内肾脏可恢复\omterm{def:b2:blood-circulation:blood}{血浆}容量。
        指出负责此作用的两种激素及其触发因素。
  \item 若骨髓以
        每秒 250 万个（\num{5e12} 每升）的速率补充\omterm{def:b2:blood-circulation:blood}{红细胞}，需多长时间？
  \item 失血 \qty{2.5}{L} 使反射饱和。结合
        增益与反射曲线，解释压强为何随之骤降。
\end{enumerate}

\medskip
\noindent\textbf{第四部分——跑步。}
在剧烈运动中，肌肉阻力降至 \qty{0.33}{mmHg.s/mL}，
皮肤阻力降至 5，肠道与肾脏阻力加倍，脑部与\omterm{def:b2:heart:anatomy}{心脏}阻力下降，
使脑部血流不变，心肌血流
在新压强下增至四倍；$\bar P$ 升至 \qty{110}{mmHg}。
\begin{enumerate}[resume]
  \item 计算流向肌肉与皮肤的血流。
  \item 计算总阻力与\omterm{thm:b2:heart:starling}{心输出量}。
  \item 心率为 180 时，计算\omterm{prop:b2:heart:cycle}{每搏输出量}。
  \item \omterm{def:b2:blood-circulation:vessels}{静脉}血氧降至 \qty{40}{mL/L}，计算摄氧量
        及其超出静息水平的倍数。
  \item 将该倍数归因于血流与提取。
  \item 解释\omterm{def:b2:blood-pressure:baroreflex}{压力反射}如何能够允许压强为
        \qty{110}{mmHg}，而非将其纠正。
  \item 归纳结果：站起时的剩余下降量、
        \qty{1}{L} 失血后的压强，以及
        运动时的\omterm{thm:b2:heart:starling}{心输出量}与摄氧量。
\end{enumerate}
\end{problem}
